\documentclass[10pt,a4paper]{amsart}
\usepackage{amsmath,amssymb,amsthm,graphicx}
\usepackage{verbatim}
\newenvironment{reference}{\comment}{\endcomment}
%\newenvironment{reference}{}{}
\newenvironment{slogan}{\comment}{\endcomment}
\newenvironment{history}{\comment}{\endcomment}
% For commutative diagrams we use Xy-pic
\usepackage[all]{xy}

% We use 2cell for 2-commutative diagrams.
\xyoption{2cell}
% Theorem environments.
%
\theoremstyle{plain}
\newtheorem{theorem}[subsection]{Theorem}
\newtheorem{proposition}[subsection]{Proposition}
\newtheorem{lemma}[subsection]{Lemma}

\theoremstyle{definition}
\newtheorem{definition}[subsection]{Definition}
\newtheorem{example}[subsection]{Example}
\newtheorem{exercise}[subsection]{Exercise}
\newtheorem{situation}[subsection]{Situation}

\theoremstyle{remark}
\newtheorem{remark}[subsection]{Remark}
\newtheorem{remarks}[subsection]{Remarks}

\numberwithin{equation}{subsection}

% Macros
%
\def\lim{\mathop{\mathrm{lim}}\nolimits}
\def\colim{\mathop{\mathrm{colim}}\nolimits}
\def\Spec{\mathop{\mathrm{Spec}}}
\def\Hom{\mathop{\mathrm{Hom}}\nolimits}
\def\Ext{\mathop{\mathrm{Ext}}\nolimits}
\def\SheafHom{\mathop{\mathcal{H}\!\mathit{om}}\nolimits}
\def\SheafExt{\mathop{\mathcal{E}\!\mathit{xt}}\nolimits}
\def\Sch{\mathit{Sch}}
\def\Mor{\mathop{\mathrm{Mor}}\nolimits}
\def\Ob{\mathop{\mathrm{Ob}}\nolimits}
\def\Sh{\mathop{\mathit{Sh}}\nolimits}
\def\NL{\mathop{N\!L}\nolimits}
\def\CH{\mathop{\mathrm{CH}}\nolimits}
\def\proetale{{pro\text{-}\acute{e}tale}}
\def\etale{{\acute{e}tale}}
\def\QCoh{\mathit{QCoh}}
\def\Ker{\mathop{\mathrm{Ker}}}
\def\Im{\mathop{\mathrm{Im}}}
\def\Coker{\mathop{\mathrm{Coker}}}
\def\Coim{\mathop{\mathrm{Coim}}}

% Boxtimes
%
\DeclareMathSymbol{\boxtimes}{\mathbin}{AMSa}{"02}

%
% Macros for moduli stacks/spaces
%
\def\QCohstack{\mathcal{QC}\!\mathit{oh}}
\def\Cohstack{\mathcal{C}\!\mathit{oh}}
\def\Spacesstack{\mathcal{S}\!\mathit{paces}}
\def\Quotfunctor{\mathrm{Quot}}
\def\Hilbfunctor{\mathrm{Hilb}}
\def\Curvesstack{\mathcal{C}\!\mathit{urves}}
\def\Polarizedstack{\mathcal{P}\!\mathit{olarized}}
\def\Complexesstack{\mathcal{C}\!\mathit{omplexes}}
% \Pic is the operator that assigns to X its picard group, usage \Pic(X)
% \Picardstack_{X/B} denotes the Picard stack of X over B
% \Picardfunctor_{X/B} denotes the Picard functor of X over B
\def\Pic{\mathop{\mathrm{Pic}}\nolimits}
\def\Picardstack{\mathcal{P}\!\mathit{ic}}
\def\Picardfunctor{\mathrm{Pic}}
\def\Deformationcategory{\mathcal{D}\!\mathit{ef}}
\usepackage{cleveref}
\crefname{originallemma}{original lemma}{original lemmas}
\Crefname{originallemma}{Original lemma}{Original lemmas}
\makeatletter
\newlabel{algebra-lemma-support-point}{{00L5}{source}{}{sourcepointer-1}{}}
\newlabel{algebra-lemma-support-point@cref}{{[originallemma][0][]00L5}{[1][0][]source}{}{}{}}
\makeatother
\crefname{originallemma}{original lemma}{original lemmas}
\Crefname{originallemma}{Original lemma}{Original lemmas}
\makeatletter
\newlabel{coherent-lemma-coherent-support-closed}{{01Y5}{source}{}{sourcepointer-2}{}}
\newlabel{coherent-lemma-coherent-support-closed@cref}{{[originallemma][0][]01Y5}{[1][0][]source}{}{}{}}
\makeatother
\crefname{originallemma}{original lemma}{original lemmas}
\Crefname{originallemma}{Original lemma}{Original lemmas}
\makeatletter
\newlabel{divisors-lemma-components-locally-finite}{{0BE1}{source}{}{sourcepointer-3}{}}
\newlabel{divisors-lemma-components-locally-finite@cref}{{[originallemma][0][]0BE1}{[1][0][]source}{}{}{}}
\makeatother
\crefname{originallemma}{original lemma}{original lemmas}
\Crefname{originallemma}{Original lemma}{Original lemmas}
\makeatletter
\newlabel{lemma-multiplicity-finite}{{02QT}{source}{}{sourcepointer-4}{}}
\newlabel{lemma-multiplicity-finite@cref}{{[originallemma][0][]02QT}{[1][0][]source}{}{}{}}
\makeatother
\crefname{originallemma}{original lemma}{original lemmas}
\Crefname{originallemma}{Original lemma}{Original lemmas}
\makeatletter
\newlabel{schemes-lemma-reduced-closed-subscheme}{{01J3}{source}{}{sourcepointer-5}{}}
\newlabel{schemes-lemma-reduced-closed-subscheme@cref}{{[originallemma][0][]01J3}{[1][0][]source}{}{}{}}
\makeatother
\crefname{originallemma}{original lemma}{original lemmas}
\Crefname{originallemma}{Original lemma}{Original lemmas}
\makeatletter
\newlabel{schemes-lemma-specialize-points}{{01J7}{source}{}{sourcepointer-6}{}}
\newlabel{schemes-lemma-specialize-points@cref}{{[originallemma][0][]01J7}{[1][0][]source}{}{}{}}
\makeatother
\crefname{originalreference}{original reference}{original references}
\Crefname{originalreference}{Original reference}{Original references}
\makeatletter
\newlabel{situation-setup}{{02QL}{source}{}{sourcepointer-7}{}}
\newlabel{situation-setup@cref}{{[originalreference][0][]02QL}{[1][0][]source}{}{}{}}
\makeatother
\title{Local finiteness of the support of a coherent module and finiteness of the length at a generic point}
\author{The Stacks Project Authors}
\date{stacks-project:chow.tex:local-snapshot-20261012; GNU FDL 1.2 or later}
\begin{document}
\maketitle
\noindent\textit{Editorial note.} All delivered mathematical text is the author's own, copied from the pinned Stacks Project chapter file chow.tex; no statement or proof step was rewritten, reordered or removed.\par\medskip
\noindent\textit{Editorial note.} Preamble provenance: the pinned chapter file opens with \textbackslash{}input\{preamble\} and carries no preamble text of its own, so the project's single Stacks-wide preamble staged at sources/stacks-chow/preamble.tex is excerpted verbatim as three spans, lines 8-12 (the verbatim package together with the project's own reference, slogan and history display-environment definitions, including the author's commented-out older form of the reference environment), lines 14-18 (Xy-pic setup) and lines 154-217 (the native theorem declarations and mathematical macros). The lines 8-12 span is required because a counted statement or its surrounding text may use those display environments, which the project typesets as comments. Nothing was added, removed or redefined; the wrapper class and the amsmath/amssymb/amsthm/graphicx base packages are supplied by the assembler.\par\medskip
\noindent\textit{Editorial note.} Same-chapter pointers are located at their \textbackslash{}label lines inside the pinned chapter file; cross-chapter pointers are located in the project's own tag list staged in the snapshot (sources/stacks-chow/tags.txt), at the verbatim lines carrying their canonical tags.\par\medskip
\noindent\textit{Allowed standard prerequisites (editorial).}\begin{itemize}
\item The support of a coherent module is closed, and its reduced structure is a reduced closed subscheme. Cohomology of Schemes, label coherent-lemma-coherent-support-closed (proof line 1723); Schemes, label schemes-lemma-reduced-closed-subscheme (proof line 1726). Realises the support as a reduced closed subscheme, so that statements about locally finite collections of components and about multiplicities apply to it.
\item Local finiteness of the irreducible components of a closed subscheme, finiteness of multiplicities, specialisation of points and finiteness of length over a ring whose module is supported at the closed point. Divisors, label divisors-lemma-components-locally-finite (proof line 1728); Same chapter chow.tex, label lemma-multiplicity-finite (proof line 1730); Schemes, label schemes-lemma-specialize-points (proof line 1737); Algebra, label algebra-lemma-support-point (proof line 1740). The first two give parts (1) and (3) directly, and the last two show that the stalk at the generic point is a finite module over a local ring supported only at the closed point, hence of finite length.
\end{itemize}\medskip
\ifdefined\hypertarget\hypertarget{sourcepointer-1}{}\fi
\noindent\textit{Source reference.} Original lemma 00L5: sources/stacks-chow/tags.txt, source lines 775--775. Navigation pointer for \textbackslash{}ref\{algebra-lemma-support-point\} used inside the counted span (line 1740); the argument that consumes it is retained in the span.\par\medskip
\ifdefined\hypertarget\hypertarget{sourcepointer-2}{}\fi
\noindent\textit{Source reference.} Original lemma 01Y5: sources/stacks-chow/tags.txt, source lines 2544--2544. Navigation pointer for \textbackslash{}ref\{coherent-lemma-coherent-support-closed\} used inside the counted span (line 1723); the argument that consumes it is retained in the span.\par\medskip
\ifdefined\hypertarget\hypertarget{sourcepointer-3}{}\fi
\noindent\textit{Source reference.} Original lemma 0BE1: sources/stacks-chow/tags.txt, source lines 14334--14334. Navigation pointer for \textbackslash{}ref\{divisors-lemma-components-locally-finite\} used inside the counted span (line 1728); the argument that consumes it is retained in the span.\par\medskip
\ifdefined\hypertarget\hypertarget{sourcepointer-4}{}\fi
\noindent\textit{Source reference.} Original lemma 02QT: sources/stacks-chow/pinned-original-source.tex, source lines 1629--1629. Navigation pointer for \textbackslash{}ref\{lemma-multiplicity-finite\} used inside the counted span (line 1730); the argument that consumes it is retained in the span.\par\medskip
\ifdefined\hypertarget\hypertarget{sourcepointer-5}{}\fi
\noindent\textit{Source reference.} Original lemma 01J3: sources/stacks-chow/tags.txt, source lines 2000--2000. Navigation pointer for \textbackslash{}ref\{schemes-lemma-reduced-closed-subscheme\} used inside the counted span (line 1726); the argument that consumes it is retained in the span.\par\medskip
\ifdefined\hypertarget\hypertarget{sourcepointer-6}{}\fi
\noindent\textit{Source reference.} Original lemma 01J7: sources/stacks-chow/tags.txt, source lines 2004--2004. Navigation pointer for \textbackslash{}ref\{schemes-lemma-specialize-points\} used inside the counted span (line 1737); the argument that consumes it is retained in the span.\par\medskip
\ifdefined\hypertarget\hypertarget{sourcepointer-7}{}\fi
\noindent\textit{Source reference.} Original reference 02QL: sources/stacks-chow/pinned-original-source.tex, source lines 1372--1372. Navigation pointer for \textbackslash{}ref\{situation-setup\} used inside the counted span (line 1703); the argument that consumes it is retained in the span.\par\medskip
\begin{lemma}
\label{lemma-length-finite}
Let $(S, \delta)$ be as in Situation \ref{situation-setup}.
Let $X$ be locally of finite type over $S$.
Let $\mathcal{F}$ be a coherent $\mathcal{O}_X$-module.
\begin{enumerate}
\item The collection of irreducible components of the support of
$\mathcal{F}$ is locally finite.
\item Let $Z' \subset \text{Supp}(\mathcal{F})$
be an irreducible component and
let $\xi \in Z'$ be its generic point.
Then
$$
\text{length}_{\mathcal{O}_{X, \xi}} \mathcal{F}_\xi < \infty
$$
\item If $\dim_\delta(\text{Supp}(\mathcal{F})) \leq k$
and $\xi \in \text{Supp}(\mathcal{F})$ with $\delta(\xi) = k$, then $\xi$ is a
generic point of an irreducible component of $\text{Supp}(\mathcal{F})$.
\end{enumerate}
\end{lemma}

\begin{proof}
By Cohomology of Schemes, Lemma \ref{coherent-lemma-coherent-support-closed}
the support $Z$ of $\mathcal{F}$ is a closed subset of $X$.
We may think of $Z$ as a reduced closed subscheme of $X$
(Schemes, Lemma \ref{schemes-lemma-reduced-closed-subscheme}).
Hence (1) follows from
Divisors, Lemma \ref{divisors-lemma-components-locally-finite} applied to $Z$
and (3) follows from
Lemma \ref{lemma-multiplicity-finite} applied to $Z$.

\medskip\noindent
Let $\xi \in Z'$ be as in (2). In this case for any specialization
$\xi' \leadsto \xi$ in $X$ we have $\mathcal{F}_{\xi'} = 0$.
Recall that the non-maximal primes of $\mathcal{O}_{X, \xi}$ correspond
to the points of $X$ specializing to $\xi$
(Schemes, Lemma \ref{schemes-lemma-specialize-points}).
Hence $\mathcal{F}_\xi$ is a finite $\mathcal{O}_{X, \xi}$-module
whose support is $\{\mathfrak m_\xi\}$. Hence it has finite length
by Algebra, Lemma \ref{algebra-lemma-support-point}.
\end{proof}

\end{document}
